\documentclass[11pt, a4paper]{report}

% ========== Common Packages ==========
\usepackage{fontspec}
\usepackage{kotex}
\usepackage{amsmath, amssymb, amsthm}
\usepackage{mathtools}
\usepackage{bm}
\usepackage{graphicx}
\usepackage{booktabs}
\usepackage{array}
\usepackage{tabularx}
\usepackage{longtable}
\usepackage{hyperref}
\usepackage{cleveref}
\usepackage{geometry}
\usepackage{fancyhdr}
\usepackage{titlesec}
\usepackage{enumitem}
\usepackage{xcolor}
\usepackage{tcolorbox}
\usepackage{algorithm}
\usepackage{algorithmic}
\usepackage{tikz}
\usetikzlibrary{arrows.meta, positioning, calc, decorations.pathreplacing, shapes.geometric, fit, patterns, angles, quotes, trees}
\usepackage{pgfplots}
\pgfplotsset{compat=1.18}
\usepackage{caption}
\usepackage{subcaption}
\usepackage{float}

\geometry{margin=2.5cm, top=3cm, bottom=3cm}

\usepackage{multirow}
% ========== Colors ==========
% Common
\definecolor{defblue}{RGB}{30, 80, 150}
\definecolor{thmgreen}{RGB}{20, 110, 60}
\definecolor{noteorange}{RGB}{200, 100, 0}
\definecolor{lightblue}{RGB}{230, 240, 255}
\definecolor{lightgreen}{RGB}{230, 255, 235}
\definecolor{lightorange}{RGB}{255, 245, 230}
\definecolor{lightgray}{RGB}{245, 245, 245}
% Extended
\definecolor{lightred}{RGB}{255, 235, 235}
\definecolor{darkred}{RGB}{180, 30, 30}
\definecolor{biogreen}{RGB}{10, 130, 80}
\definecolor{cvpurple}{RGB}{100, 40, 150}
\definecolor{injuryred}{RGB}{200, 40, 40}
\definecolor{codegray}{RGB}{240, 240, 245}
\definecolor{pipeblue}{RGB}{25, 100, 180}
\definecolor{constraintred}{RGB}{200, 50, 50}
\definecolor{termblack}{RGB}{30, 30, 30}
\definecolor{goldcolor}{RGB}{180, 140, 20}

% ========== tcolorbox environments ==========
% --- Common (all parts) ---
\newtcolorbox{keyidea}[1][]{
  colback=lightblue,
  colframe=defblue,
  fonttitle=\bfseries,
  title={Key Idea},
  #1
}

\newtcolorbox{importantnote}[1][]{
  colback=lightorange,
  colframe=noteorange,
  fonttitle=\bfseries,
  title={Important},
  #1
}

\newtcolorbox{summary}[1][]{
  colback=lightgreen,
  colframe=thmgreen,
  fonttitle=\bfseries,
  title={Summary},
  #1
}

% --- Warning/Error ---
\newtcolorbox{warningbox}[1][]{
  colback=lightred,
  colframe=darkred,
  fonttitle=\bfseries,
  title={Warning},
  #1
}

% --- Domain: Biomechanics & Clinical ---
\newtcolorbox{clinicalnote}[1][]{
  colback=lightred,
  colframe=darkred,
  fonttitle=\bfseries,
  title={Clinical Note},
  #1
}

\newtcolorbox{biomechbox}[1][]{
  colback=lightgreen,
  colframe=biogreen,
  fonttitle=\bfseries,
  title={Biomechanics},
  #1
}

\newtcolorbox{injurybox}[1][]{
  colback={injuryred!8},
  colframe=injuryred,
  fonttitle=\bfseries,
  title={Injury Risk},
  #1
}

% --- Domain: Computer Vision ---
\newtcolorbox{cvbox}[1][]{
  colback={cvpurple!5},
  colframe=cvpurple,
  fonttitle=\bfseries,
  title={Computer Vision},
  #1
}

% --- Pipeline & Setup ---
\newtcolorbox{setupbox}[1][]{
  colback=lightgray,
  colframe=gray!60,
  fonttitle=\bfseries,
  title={Setup},
  #1
}

\newtcolorbox{pipelinebox}[1][]{
  colback=pipeblue!5,
  colframe=pipeblue,
  fonttitle=\bfseries,
  title={Pipeline Step},
  #1
}

\newtcolorbox{codebox}[1][]{
  colback=codegray,
  colframe=gray!70,
  fonttitle=\bfseries\ttfamily,
  title={Code},
  #1
}

\newtcolorbox{termbox}[1][]{
  colback=termblack!5,
  colframe=termblack!60,
  fonttitle=\bfseries\ttfamily,
  title={Terminal},
  #1
}

% --- Research ---
\newtcolorbox{hypothesis}[1][]{
  colback=goldcolor!8,
  colframe=goldcolor,
  fonttitle=\bfseries,
  title={Hypothesis},
  #1
}

\newtcolorbox{resultbox}[1][]{
  colback=thmgreen!8,
  colframe=thmgreen,
  fonttitle=\bfseries,
  title={Expected Result},
  #1
}

% --- Constraints & Solutions ---
\newtcolorbox{constraintbox}[1][]{
  colback=constraintred!8,
  colframe=constraintred,
  fonttitle=\bfseries,
  title={Constraint},
  #1
}

\newtcolorbox{solutionbox}[1][]{
  colback=thmgreen!8,
  colframe=thmgreen,
  fonttitle=\bfseries,
  title={Solution},
  #1
}

\newtcolorbox{databox}[1][]{
  colback=pipeblue!5,
  colframe=pipeblue,
  fonttitle=\bfseries,
  title={Data Spec},
  #1
}

% --- Progress/Status ---
\newtcolorbox{successbox}[1][]{
  colback=lightgreen,
  colframe=thmgreen,
  fonttitle=\bfseries,
  title={Success},
  #1
}

\newtcolorbox{nextbox}[1][]{
  colback=lightorange,
  colframe=noteorange,
  fonttitle=\bfseries,
  title={Next Step},
  #1
}

% --- Fitting guide ---
\newtcolorbox{failbox}[1][]{
  colback=lightred,
  colframe=darkred,
  fonttitle=\bfseries,
  title={Failure Pattern},
  #1
}

\newtcolorbox{fixbox}[1][]{
  colback=lightgreen,
  colframe=thmgreen,
  fonttitle=\bfseries,
  title={Fix},
  #1
}

\newtcolorbox{lessonbox}[1][]{
  colback=lightorange,
  colframe=noteorange,
  fonttitle=\bfseries,
  title={Lesson Learned},
  #1
}

% --- Specification ---
\newtcolorbox{specbox}[1][]{
  colback=cvpurple!5,
  colframe=cvpurple,
  fonttitle=\bfseries,
  title={Specification},
  #1
}

\newtcolorbox{checklistbox}[1][]{
  colback=lightgreen,
  colframe=biogreen,
  fonttitle=\bfseries,
  title={Checklist},
  #1
}

% ========== Custom math commands ==========
% Vectors (lowercase)
\newcommand{\vx}{\mathbf{x}}
\newcommand{\vd}{\mathbf{d}}
\newcommand{\vo}{\mathbf{o}}
\newcommand{\vc}{\mathbf{c}}
\newcommand{\vr}{\mathbf{r}}
\newcommand{\vp}{\mathbf{p}}
\newcommand{\vv}{\mathbf{v}}
\newcommand{\vn}{\mathbf{n}}
\newcommand{\vu}{\mathbf{u}}
\newcommand{\vt}{\mathbf{t}}
\newcommand{\ve}{\mathbf{e}}
\newcommand{\vq}{\mathbf{q}}
% Vectors (uppercase)
\newcommand{\vF}{\mathbf{F}}
\newcommand{\vM}{\mathbf{M}}
\newcommand{\va}{\mathbf{a}}
\newcommand{\vg}{\mathbf{g}}
\newcommand{\vX}{\mathbf{X}}
% Bold Greek
\newcommand{\vmu}{\bm{\mu}}
\newcommand{\vbeta}{\bm{\beta}}
\newcommand{\vtheta}{\bm{\theta}}
\newcommand{\vpsi}{\bm{\psi}}
% Matrices
\newcommand{\mSigma}{\bm{\Sigma}}
\newcommand{\mR}{\mathbf{R}}
\newcommand{\mS}{\mathbf{S}}
\newcommand{\mW}{\mathbf{W}}
\newcommand{\mJ}{\mathbf{J}}
\newcommand{\mG}{\mathbf{G}}
\newcommand{\mT}{\mathbf{T}}
\newcommand{\mI}{\mathbf{I}}
\newcommand{\mB}{\mathbf{B}}
\newcommand{\mK}{\mathbf{K}}
\newcommand{\mP}{\mathbf{P}}
\newcommand{\mF}{\mathbf{F}}
\newcommand{\mE}{\mathbf{E}}
\newcommand{\mA}{\mathbf{A}}
\newcommand{\mH}{\mathbf{H}}
% Sets and groups
\newcommand{\RR}{\mathbb{R}}
\newcommand{\SO}{\mathrm{SO}}
% Units
\newcommand{\degps}{\,\text{deg/s}}
\newcommand{\Nm}{\ensuremath{\,\text{N}\cdot\text{m}}}
\newcommand{\BW}{\,\text{BW}}

% ========== Listings ==========
\usepackage{listings}

\lstset{
  basicstyle=\ttfamily\small,
  keywordstyle=\color{defblue}\bfseries,
  commentstyle=\color{thmgreen}\itshape,
  stringstyle=\color{noteorange},
  backgroundcolor=\color{codegray},
  frame=single,
  rulecolor=\color{gray!40},
  breaklines=true,
  showstringspaces=false,
  numbers=left,
  numberstyle=\tiny\color{gray},
  tabsize=4,
  captionpos=b
}

% ========== Header/Footer ==========
\pagestyle{fancy}
\fancyhf{}
\fancyhead[L]{\leftmark}
\fancyhead[R]{\thepage}
\renewcommand{\headrulewidth}{0.4pt}

% ========== Title formatting ==========
\titleformat{\chapter}[display]
  {\normalfont\huge\bfseries\color{defblue}}
  {\chaptertitlename\ \thechapter}{20pt}{\Huge}

\titleformat{\section}
  {\normalfont\Large\bfseries\color{defblue!80!black}}
  {\thesection}{1em}{}

\titleformat{\subsection}
  {\normalfont\large\bfseries\color{defblue!60!black}}
  {\thesubsection}{1em}{}


\begin{document}

% ==================== TITLE PAGE ====================
\begin{titlepage}
\centering
\vspace*{1.5cm}
{\Large\color{gray} Practical Fitting Guide}
\vspace{0.5cm}

{\Huge\bfseries\color{defblue} SMPL 피팅 실전 가이드\\[0.3cm]
\&{} 오답노트}

\vspace{0.8cm}

{\LARGE\color{defblue!70!black} 시행착오에서 배운 것들:\\
무엇이 실패했고, 왜 실패했고, 어떻게 해야 하는가}

\vspace{2cm}

\begin{tikzpicture}[
  step/.style={rectangle, rounded corners=4pt, draw=defblue, fill=lightblue,
    minimum width=2cm, minimum height=0.6cm, font=\scriptsize\bfseries, align=center},
  fail/.style={circle, draw=darkred, fill=lightred, minimum size=0.4cm, font=\tiny\bfseries, text=darkred},
  arrow/.style={-Stealth, thick, defblue!60}
]
\node[step] (cal) at (0,0) {카메라\\보정};
\node[step] (pose) at (2.5,0) {2D 자세\\추정};
\node[step] (tri) at (5,0) {3D\\삼각측량};
\node[step] (smpl) at (7.5,0) {SMPL\\피팅};
\node[step] (gart) at (10,0) {3DGS\\정제};
\node[step] (bio) at (12.5,0) {생체역학\\분석};

\draw[arrow] (cal) -- (pose);
\draw[arrow] (pose) -- (tri);
\draw[arrow] (tri) -- (smpl);
\draw[arrow] (smpl) -- (gart);
\draw[arrow] (gart) -- (bio);

\node[fail] at (0, 0.7) {!};
\node[fail] at (2.5, 0.7) {!};
\node[fail] at (5, 0.7) {!};
\node[fail] at (7.5, 0.7) {!};
\node[fail] at (10, 0.7) {!};

\node[font=\tiny\color{darkred}] at (0, 1.1) {COLMAP 실패};
\node[font=\tiny\color{darkred}] at (2.5, 1.1) {인물 선택 오류};
\node[font=\tiny\color{darkred}] at (5, 1.1) {좌표계 불일치};
\node[font=\tiny\color{darkred}] at (7.5, 1.1) {FK 오류};
\node[font=\tiny\color{darkred}] at (10, 1.1) {미진행};
\end{tikzpicture}

\vspace{2cm}
{\large 2026년 4월}

\vspace{\fill}
{\small\color{gray} 실제 데이터(7대 240Hz + Vicon)에 파이프라인을 적용하면서\\
겪은 모든 실패와 그로부터 배운 교훈을 기록합니다.}
\end{titlepage}

\tableofcontents
\newpage


% ====================================================================
\chapter{피팅 전략 총론}
% ====================================================================

\section{교재에서 다룬 피팅 방법 계보}

\begin{table}[H]
\centering
\caption{SMPL 피팅 방법 계보와 특성}
\renewcommand{\arraystretch}{1.3}
\small
\begin{tabularx}{\textwidth}{lXccl}
\toprule
\textbf{방법} & \textbf{핵심 아이디어} & \textbf{입력} & \textbf{출력} & \textbf{교재} \\
\midrule
SMPLify & 2D kp $\rightarrow$ 재투영 최적화 & 2D & $\vtheta, \vbeta$ & Part 2 \\
HMR/SPIN & End-to-end 회귀 & 이미지 & $\vtheta, \vbeta$ & Part 2 \\
EasyMocap & 다시점 2D $\rightarrow$ SMPL-X & 2D $\times$ C & $\vtheta, \vbeta$ & Part 3 \\
\textbf{3D 직접 피팅} & \textbf{3D kp $\rightarrow$ SMPL 최적화} & \textbf{3D} & $\vtheta, \vbeta$ & \textbf{Part 4} \\
3DGS 자세 정제 & Photometric loss로 정제 & 영상 & $\vtheta^*$ & Part 5 \\
\bottomrule
\end{tabularx}
\end{table}

\section{현재 세팅에서 최적 경로}

\begin{keyidea}
\textbf{현재 보유한 것}:
\begin{itemize}[nosep]
  \item 카메라 보정: DA3 PnP refined (검증됨)
  \item 세그멘테이션: SAM3 마스크 7대 $\times$ 1000프레임 (완료)
  \item 2D 자세: OpenPose Body25 $\times$ 7대 (완료, cam7 tracked)
  \item 3D 자세: DLT 삼각측량 25관절 (검증됨)
  \item SMPL 파라미터: USB에 여러 버전 존재 (미검증)
  \item \textbf{누락}: smplx 라이브러리가 로컬에 없어 FK 검증 불가
\end{itemize}

\textbf{최적 경로}:
\begin{enumerate}[nosep]
  \item 서버(A100)에서 smplx로 기존 파라미터의 FK를 실행하여 관절 위치 추출
  \item OpenPose 3D와 비교하여 올바른 파라미터 식별
  \item 필요시 서버에서 3D 직접 피팅 재실행 (smplx + PyTorch)
\end{enumerate}
\end{keyidea}


% ====================================================================
\chapter{시행착오 기록 (오답노트)}
% ====================================================================

\section{카메라 보정}

\begin{failbox}[title={실패 1: COLMAP SfM 카메라 보정}]
\textbf{시도}: COLMAP으로 7대 카메라의 intrinsic/extrinsic 추정\\
\textbf{결과}: 등록률 100\% 달성했으나, GT(DA3 PnP)와 비교 시 카메라 위치 오차 50\%+, 초점거리 47--62\% 오차\\
\textbf{원인}:
\begin{enumerate}[nosep]
  \item 동적 장면(투수 움직임)이 SfM의 정적 장면 가정 위반
  \item 7대 카메라가 동일 기종인데 COLMAP이 인식 못해 초점거리를 531--1367px로 추정 (GT: $\sim$790px 전체 일관)
  \item Round 1 (OPENCV 8파라미터)에서 과적합 $\rightarrow$ $f_x \neq f_y$
\end{enumerate}
\end{failbox}

\begin{fixbox}
\textbf{해결}: DA3 PnP refined 결과를 채택.\\
\textbf{교훈}: 동적 장면에서 COLMAP SfM은 부적합. 기존 보정 결과가 있으면 그것을 사용.\\
\textbf{추가 작업}: 280$\times$280 $\rightarrow$ 원본 해상도 역변환 (upper\_bound\_resize + center\_crop 추적)
\end{fixbox}


\section{세그멘테이션}

\begin{failbox}[title={실패 2: cam7 SAM3 마스크 방향 불일치}]
\textbf{시도}: 기존 SAM3 마스크를 그대로 사용\\
\textbf{결과}: cam7 마스크가 portrait(1080$\times$1920)로 생성되어 이미지(landscape 1920$\times$1080)와 방향 불일치\\
\textbf{추가 문제}: 투수가 아닌 다른 사람이 간헐적으로 선택됨 (cam7 시야에 3명)
\end{failbox}

\begin{fixbox}
\textbf{해결}: Tracked OpenPose(P0=투수) + SAM3를 서버에서 재실행 $\rightarrow$ 1000/1000 성공, 올바른 해상도\\
\textbf{교훈}: 다인원 장면에서 person selection이 핵심. OpenPose tracking + SAM3 text prompt ``person'' + 키포인트 겹침 선택.
\end{fixbox}


\section{2D 자세 추정}

\begin{failbox}[title={실패 3: ViTPose person 선택 실패}]
\textbf{시도}: ViTPose(MMPose)로 원본 이미지에서 2D 키포인트 추출\\
\textbf{결과}: cam3에서 90.9\%의 프레임에서 투수가 아닌 다른 사람이 P0으로 선택됨\\
\textbf{원인}: ViTPose가 detection confidence 순으로 정렬하는데, 정적인 다른 사람의 confidence가 더 높음
\end{failbox}

\begin{fixbox}
\textbf{시도한 해결}: masked images (배경 제거)에서 ViTPose 재실행\\
\textbf{결과}: 투수만 검출되지만, OpenPose 3D 대비 jitter 2배, sudden jump 30배 $\rightarrow$ 품질 열등\\
\textbf{교훈}: 기존 OpenPose + pitcher\_indices가 더 안정적. ``더 좋은 모델''이 항상 더 좋은 결과를 주지 않음.\\
\textbf{참고}: MMPose 기본 모델은 RTMPose이지 ViTPose-Huge가 아니었음.
\end{fixbox}


\section{SMPL 피팅 --- 핵심 난관}

\begin{failbox}[title={실패 4: EasyMocap SMPL 좌표계 불일치}]
\textbf{데이터}: CJH4/easymocap\_smplx\_init.npz (998프레임)\\
\textbf{결과}: OpenPose 3D MidHip $(0.12, 2.81, 1.21)$ vs SMPL Th $(0.16, 0.79, 2.71)$ $\rightarrow$ \textbf{Y축과 Z축이 뒤바뀜}\\
\textbf{원인}: EasyMocap이 2D에서 3D를 재구성하면서 다른 좌표 규약 사용. 글로벌 회전(Rh) 평균 123$^\circ$로 좌표계 보정을 흡수.
\end{failbox}

\begin{failbox}[title={실패 5: Manual FK로 SMPL 관절 위치 계산 실패}]
\textbf{시도}: smplx 라이브러리 없이 수동 Forward Kinematics 구현 (REST pose 뼈 길이를 임의 근사)\\
\textbf{결과}: CJH2/optimized\_params, CJH3/vicon\_solved, CJH4/easymocap 모두 3D 뷰어에서 비현실적 스켈레톤\\
\textbf{원인}:
\begin{enumerate}[nosep]
  \item REST pose의 뼈 길이가 \textbf{임의의 근사값}이라 실제 SMPL 토폴로지와 불일치
  \item SMPL의 blend shapes ($B_S(\vbeta), B_P(\vtheta)$)를 무시하여 vertex deformation 미반영
  \item 관절 위치가 regressor ($J = \mathcal{J}\bar{T}$)로 결정되는데, 단순 FK 체인으로는 정확한 위치 불가능
\end{enumerate}
\end{failbox}

\begin{failbox}[title={실패 6: 새 SMPL 피팅 (서버) --- MPJPE 219.6mm}]
\textbf{시도}: OpenPose 3D $\rightarrow$ smplx + PyTorch로 직접 피팅 (4-stage coarse-to-fine)\\
\textbf{결과}: MPJPE 219.6mm (발목 391--442mm, 무릎 293--300mm)\\
\textbf{원인 추정}:
\begin{enumerate}[nosep]
  \item SMPL 모델이 SMPL\_NEUTRAL.pkl (304KB)로, 전체 모델이 아닌 경량 버전일 가능성
  \item $\vbeta$ 최적화가 불충분 (뼈 길이가 피험자 체형과 불일치)
  \item Iteration 수 부족 (1400 total)
  \item 자세 사전($E_\theta$)이 투구 동작에 부적합
\end{enumerate}
\end{failbox}

\begin{lessonbox}
\textbf{핵심 교훈: smplx 라이브러리의 Forward Kinematics가 반드시 필요하다}

SMPL 파라미터 $(\vtheta, \vbeta, \mR, \vt)$에서 관절 위치를 정확히 계산하려면
\textbf{smplx 라이브러리의 forward 함수를 사용해야 한다}:

\begin{lstlisting}[language=Python, numbers=none]
output = body_model(
    global_orient=global_orient,  # (N, 3)
    body_pose=body_pose,          # (N, 63) or (N, 69)
    betas=betas,                  # (1, 10)
    transl=transl                 # (N, 3)
)
joints = output.joints[:, :24, :]  # Correct 3D positions
\end{lstlisting}

수동 FK는 blend shapes, joint regressor, pose-dependent corrections를 무시하므로
\textbf{절대 정확한 결과를 줄 수 없다}.
\end{lessonbox}


% ====================================================================
\chapter{올바른 피팅 전략}
% ====================================================================

\section{USB에 존재하는 SMPL 파라미터 파일 목록}

\begin{table}[H]
\centering
\caption{T31/CJH\_motioncapture에서 발견된 SMPL 파라미터}
\renewcommand{\arraystretch}{1.3}
\small
\begin{tabularx}{\textwidth}{clcXl}
\toprule
\textbf{\#} & \textbf{위치} & \textbf{프레임} & \textbf{파라미터} & \textbf{출처} \\
\midrule
1 & CJH2/optimized\_params/ & 125 & body\_pose(63), hands, jaw, expression & SMPL-X full \\
2 & CJH2/output\_gpu/params/ & \textbf{998} & body\_pose(63), hands, jaw, expression & \textbf{GPU 최적화} \\
3 & CJH2/output\_gpu\_v3/params/ & \textbf{998} & 동일 & GPU v3 \\
4 & CJH3/vicon\_solved\_v13 & 875 & body\_pose(63), betas(10) & Vicon 기반 \\
5 & CJH3/data/vicon/ (v2--v11) & 875 & 동일 & 여러 버전 \\
6 & CJH4/easymocap\_init & 998 & poses(72), Rh, Th & EasyMocap \\
7 & 4DGS/smplx\_sequence & 125 & body\_pose(63), joints(22,3) & 다시점 \\
8 & \textbf{새 피팅 (서버)} & \textbf{1000} & poses(72), betas(10) & OP3D 직접 \\
\bottomrule
\end{tabularx}
\end{table}


\section{올바른 검증 방법}

\begin{keyidea}
\textbf{모든 파라미터 파일을 GPU 서버에서 smplx FK로 관절 위치를 추출}한 뒤,
\textbf{OpenPose 3D와 비교}하여 어느 것이 가장 정확한지 판별해야 한다.

\begin{enumerate}[nosep]
  \item 서버에 파라미터 파일 업로드
  \item \texttt{smplx.create() $\rightarrow$ output.joints[:,:24,:]} 로 관절 위치 추출
  \item OpenPose 3D와 MPJPE 계산
  \item 3D 뷰어로 시각적 확인
  \item 가장 낮은 MPJPE를 가진 파일 채택
\end{enumerate}
\end{keyidea}

\begin{warningbox}
\textbf{로컬에서 수동 FK를 쓰지 말 것}. 반드시 서버의 smplx 라이브러리를 사용할 것.\\
수동 FK의 REST pose 뼈 길이는 실제 SMPL 토폴로지와 다르며,
blend shapes와 joint regressor를 무시하므로 관절 위치가 완전히 잘못된다.
\end{warningbox}


\section{피팅이 안 맞을 경우: 재피팅 전략}

기존 파라미터가 모두 부적합할 경우, 서버에서 재피팅:

\begin{enumerate}[leftmargin=*]
  \item \textbf{SMPL-X 모델 확인}: \texttt{SMPLX\_NEUTRAL.npz} 사용 (SMPL\_NEUTRAL.pkl이 아닌 풀 모델)
  \item \textbf{3D keypoints 직접 피팅} (Part 4 방법):
  \begin{equation}
  \min_{\vtheta, \vbeta, \mR, \vt} \sum_{j \in \mathcal{M}} w_j \| J_j^{\text{SMPL}}(\vtheta, \vbeta) + \vt - \mathbf{p}_j^{\text{OP3D}} \|^2 + \lambda_\beta \|\vbeta\|^2 + \lambda_\theta \|\vtheta\|^2
  \end{equation}
  \item \textbf{$\vbeta$ 충분히 최적화}: 뼈 길이가 피험자와 맞을 때까지 (1000+ iterations)
  \item \textbf{시간 정규화 적용}: 시퀀스 전체에 가속도 정규화
  \item \textbf{검증}: smplx FK로 관절 위치 추출 $\rightarrow$ OpenPose 3D와 비교 (MPJPE $< 50$mm 목표)
\end{enumerate}


% ====================================================================
\chapter{하이퍼파라미터 튜닝 가이드}
% ====================================================================

\section{SMPL 피팅 하이퍼파라미터}

\begin{table}[H]
\centering
\caption{SMPL 3D 피팅 핵심 하이퍼파라미터}
\renewcommand{\arraystretch}{1.3}
\begin{tabularx}{\textwidth}{lcXX}
\toprule
\textbf{파라미터} & \textbf{기본값} & \textbf{효과} & \textbf{투구 조정} \\
\midrule
$\lambda_\theta$ & 0.001 & 자세 정규화 (T-pose 방향) & \textbf{0.0001} (극단적 자세 허용) \\
$\lambda_\beta$ & 0.01 & 체형 정규화 (평균 방향) & 0.01 유지 \\
$\lambda_{\text{smooth}}$ & 0.1 & 시간적 매끄러움 & 가속 단계: \textbf{0.01} \\
$\gamma$ (가속도) & 0.5 & 가속도 정규화 강도 & 가속 단계: \textbf{0.1} \\
학습률 (Stage 3) & $10^{-3}$ & $\vtheta$ 최적화 속도 & $10^{-3}$ 유지 \\
학습률 (Stage 4) & $10^{-4}$ & 공동 정제 속도 & $5 \times 10^{-5}$ \\
Iterations (total) & 1400 & 수렴 충분성 & \textbf{3000+} \\
배치 크기 & 50 & GPU 메모리 사용 & A100: 100 가능 \\
\bottomrule
\end{tabularx}
\end{table}

\begin{importantnote}
\textbf{투구 동작의 가속 단계(30ms)에서는 $\lambda_\theta$와 $\lambda_{\text{smooth}}$를 줄여야 한다.}
7,000$^\circ$/s의 어깨 내회전을 정규화가 억제하면 정확도가 크게 저하된다.
위상별 적응형 가중치(Part 5 제안)가 이상적이지만,
단순하게 전체 시퀀스에서 $\lambda_\theta = 0.0001$로 설정해도 효과적이다.
\end{importantnote}


\section{SMPL 모델 선택}

\begin{table}[H]
\centering
\caption{SMPL 모델 파일 비교}
\renewcommand{\arraystretch}{1.3}
\begin{tabularx}{\textwidth}{lccX}
\toprule
\textbf{파일} & \textbf{크기} & \textbf{관절} & \textbf{비고} \\
\midrule
SMPL\_NEUTRAL.pkl & 304 KB & 24 & 경량, 일부 기능 제한 \\
SMPL\_NEUTRAL\_FIXED.pkl & 21 MB & 24 & 풀 모델, Chumpy 제거 \\
SMPLX\_NEUTRAL.npz & $\sim$50 MB & 127 & 풀 SMPL-X (손+얼굴) \\
\bottomrule
\end{tabularx}
\end{table}

\begin{warningbox}
서버의 \texttt{{\textasciitilde}/models/smpl/SMPL\_NEUTRAL.pkl}은 \textbf{304KB 경량 모델}이다.
이것으로 피팅하면 관절 위치 정확도가 떨어질 수 있다.
\texttt{EasyMocap/data/smplx/smpl/SMPL\_NEUTRAL\_FIXED.pkl} (21MB, 풀 모델)을 사용해야 한다.
\end{warningbox}


% ====================================================================
\chapter{방법별 비교}
% ====================================================================

\section{현재 세팅의 강점과 약점}

\begin{table}[H]
\centering
\caption{현재 세팅 SWOT}
\renewcommand{\arraystretch}{1.3}
\begin{tabularx}{\textwidth}{lX}
\toprule
\textbf{강점} & 7대 카메라 240Hz, Vicon 동시 촬영, DA3 PnP 보정 검증됨, OpenPose 3D 검증됨 \\
\midrule
\textbf{약점} & Vicon-RGB 시간 미동기화, 체커보드 없음, 로컬에 smplx 미설치 \\
\midrule
\textbf{기회} & DA3 점군 활용 가능, SAM3 마스크 완비, 서버(A100) 사용 가능 \\
\midrule
\textbf{위협} & SMPL FK 없이는 파라미터 검증 불가, 기존 결과의 좌표계 불일치 \\
\bottomrule
\end{tabularx}
\end{table}

\section{다음 단계 우선순위}

\begin{summary}[title={즉시 실행할 작업 (우선순위 순)}]
\begin{enumerate}[leftmargin=*]
  \item \textbf{서버에서 기존 SMPL 파라미터 검증}\\
    CJH2/output\_gpu/params (998프레임)를 서버에서 smplx FK $\rightarrow$ 관절 위치 추출 $\rightarrow$ OpenPose 3D와 MPJPE 비교.
    \textbf{이것이 맞으면 재피팅 불필요}.

  \item \textbf{맞는 파라미터가 없으면: 서버에서 재피팅}\\
    SMPL\_NEUTRAL\_FIXED.pkl (21MB) + OpenPose 3D $\rightarrow$ 3000+ iterations.

  \item \textbf{Vicon-RGB 시간 동기화}\\
    손목 궤적 크로스-상관으로 시간 오프셋 추정.

  \item \textbf{GART 학습}\\
    검증된 SMPL 파라미터 + 7대 영상 $\rightarrow$ GART.
\end{enumerate}
\end{summary}


% ====================================================================
\chapter{도구 및 코드 레퍼런스}
% ====================================================================

\section{서버 환경}

\begin{table}[H]
\centering
\renewcommand{\arraystretch}{1.3}
\begin{tabular}{ll}
\toprule
\textbf{항목} & \textbf{값} \\
\midrule
SSH & \texttt{ssh -i <pem> elicer@central-01.tcp.tunnel.elice.io -p 25341} \\
GPU & NVIDIA A100 80GB \\
Python & 3.10 \\
PyTorch & 2.5.1+cu121 \\
smplx & 설치됨 (\texttt{import smplx}) \\
MMPose & 1.3.2 (RTMPose 기본) \\
SAM3 & \texttt{{\textasciitilde}/sam3/} \\
\bottomrule
\end{tabular}
\end{table}

\section{핵심 서버 경로}

\begin{lstlisting}[language=bash, numbers=none]
# SMPL 모델
~/models/smpl/SMPL_NEUTRAL.pkl         # 304KB (경량)
~/EasyMocap/data/smplx/smpl/SMPL_NEUTRAL_FIXED.pkl  # 21MB (풀)
~/mocap/smplx/smplx/SMPLX_NEUTRAL.npz  # SMPL-X

# 데이터
~/results/1번/set_002/step01_frames/cam{1-7}/  # 원본 프레임
~/results/1번/set_002/step02_openpose/device{1-7}/  # OpenPose 2D
~/results/1번/set_002/step02_openpose_tracked/device7/  # cam7 tracked
~/CJH6_sam3/sam3_masks/cam{1-7}/  # SAM3 마스크
~/openpose3d_data.json  # OpenPose 3D (업로드됨)
~/smpl_fitting_result/  # 새 피팅 결과

# 결과물
~/vitpose_masked_run/   # ViTPose masked 결과
\end{lstlisting}


\section{SMPL FK 실행 코드 (서버용)}

\begin{lstlisting}[language=Python, caption={서버에서 SMPL 파라미터 $\rightarrow$ 관절 위치 추출}]
import smplx, torch, numpy as np

# Load model (FIXED version for accuracy)
model = smplx.create(
    '~/EasyMocap/data/smplx',
    model_type='smpl',
    gender='neutral',
    batch_size=50
).cuda()

# Load parameters
d = np.load('params.npz')
go = torch.tensor(d['global_orient'], dtype=torch.float32).cuda()
bp = torch.tensor(d['body_pose'], dtype=torch.float32).cuda()
bt = torch.tensor(d['betas'], dtype=torch.float32).cuda()
tr = torch.tensor(d['transl'], dtype=torch.float32).cuda()

# Forward kinematics
with torch.no_grad():
    out = model(global_orient=go, body_pose=bp,
                betas=bt, transl=tr)
    joints = out.joints[:, :24, :].cpu().numpy()

# Compare with OpenPose 3D
# joints shape: (N, 24, 3) -- correct 3D positions
\end{lstlisting}


\end{document}
